\documentclass[12pt]{article}
\usepackage[utf8]{inputenc}
\usepackage[T1]{fontenc}
\usepackage{hyperref}
\title{Seti Language Resources for K-12 Curriculum Development}
\author{Compiled by Claude Sonnet 5.5 (claude-sonnet-5-5) for Andr\'as Kornai}
\date{}
\begin{document}
\maketitle

\section{Introduction}
Seti (ISO 639-3 sbi, Glottocode seti1248) is a Torricelli language of Sandaun Province, Papua New Guinea. It belongs to the West Wapei (One) branch. About 160 speakers are cited (2003); Joshua Project gives 300, citing Ethnologue 2016. It is distinct from the unrelated Liana-Seti. Searches found no primers, dictionary or grammar specific to Seti.

\section{Primary Educational Resources}
No educational materials in Seti were found in search results.

\section{Linguistic and Reference Materials}
\begin{itemize}
\item Seti language (Wikipedia): \url{https://en.wikipedia.org/wiki/Seti_language}
\item One languages (Wikipedia), including Seti within West Wapei: \url{https://en.wikipedia.org/wiki/One_languages}
\item ASJP Seti entry: \url{https://asjp.clld.org/languages/SETI}
\item ASJP Seti Mauwil entry: \url{https://asjp.clld.org/languages/SETI_MAUWIL}
\end{itemize}

\section{Religious and Community Resources}
\begin{itemize}
\item Joshua Project people-group profile (Bible translation reported as still needed): \url{https://joshuaproject.net/people_groups/14876}
\item Joshua Project language page: \url{https://joshuaproject.net/languages/sbi}
\end{itemize}

\section{Digital and Audio Resources}
No audio or digital resources specific to Seti were found.

\section{Curriculum Development Context}
Not all young people speak the language, so transmission is at risk. There is no known school use. Related Wapei--Palei language documentation might guide orthography work, but materials for Seti itself would have to be created with the community.
\end{document}